% 修复实验的内容集中在这里，等 run_remedy.sh 的结果出来后只改这一个文件。

\newcommand{\REMEDYABSTRACT}{%
Both mechanisms are real, and one of them is locally actionable: reweighting the
objective towards the decision region raises enhancing-tumour Dice from $0.228$
to $0.306$, a gain that is significant in ten of twelve patient-paired cells,
but it costs $0.222$ whole-tumour Dice, and reweighting towards the amplified
late-time region degrades every region. Neither recovers the method. The
failure therefore lies in the conditioning of the regression target rather than
in the allocation of the loss.%
}

\newcommand{\REMEDYINTRO}{%
Finally we ask whether the two mechanisms are \emph{actionable}. If the collapse
were caused by the loss misallocating capacity, then correcting the allocation
should help. Section~\ref{sec:remedy-results} reweights the same objective along
the two axes the diagnostic identified and reports both what improves and what
does not.%
}

\newcommand{\REMEDYCONTRIB}{%
A test of two loss reweightings derived from those mechanisms, showing that the
decision-region mechanism is locally actionable --- it closes a fifth of the
enhancing-tumour gap to the discriminative model --- while neither reweighting
rescues whole-tumour segmentation. This locates the failure in the
parameterisation of the regression target rather than in the allocation of the
loss.%
}

\newcommand{\REMEDYSECTION}{%
\subsection{Are the diagnosed mechanisms actionable?}
\label{sec:remedy-results}

Table~\ref{tab:remedy} reports the two reweightings and their combination, each
trained at three seeds with everything else identical to the baseline. Neither
recovers whole-tumour segmentation, and the two axes behave very differently.

\paragraph{Foreground weighting is locally effective but costs the primary
metric.}
Raising the weight of non-background voxels by $10\times$ --- the direct
response to the observation that background occupies over $80\%$ of the sampled
volume at two orders of magnitude lower error --- raises enhancing-tumour Dice from
$0.228$ to $0.306$, a relative gain of $34\%$, and leaves tumour-core Dice
essentially unchanged ($0.372$ against $0.360$). Patient-paired bootstrap gives
a positive interval in ten of the twelve (seed, missing-pattern) cells for ET,
against only one negative. The mechanism is therefore real: the background did
dominate the objective, and the smallest class was the one paying for it.

But whole-tumour Dice falls from $0.553$ to $0.331$, with eleven of twelve cells
significantly negative and a mean paired difference of $-0.222$. The remedy
trades the primary metric for the rarest one; it does not rescue the method.
That the two regions move in opposite directions is itself informative: the
velocity field is a single global object, so re-weighting voxels does not
reallocate capacity within the field, it changes the field the model converges
to. The background-dominated objective was evidently not the binding constraint
on whole-tumour accuracy.

\paragraph{Weighting the amplified region makes everything worse.}
If the $1/(1-t)$ amplification of Eq.~\eqref{eq:sensitivity} were the operative problem, putting
capacity where the amplification acts should help. It does the opposite: the
late-time weight lowers whole-tumour Dice to $0.393$ (nine of twelve cells
significantly negative), tumour-core to $0.237$ (ten of twelve) and
enhancing-tumour to $0.110$ (nine of twelve). Combining both factors is worse
still on whole-tumour Dice --- $0.288$, significantly negative in all twelve
cells, mean paired difference $-0.265$ --- and returns enhancing-tumour Dice to
the baseline ($0.217$ against $0.228$). The late-time weight cancels the one
thing the foreground weight achieved.

This is the sharper result. The diagnostic established that the velocity error
grows with $t$ and that the target's sensitivity to posterior error diverges
there; the natural reading, ``train harder where it is wrong'', is wrong. The
high-$t$ target is not merely hard to fit, it is ill-conditioned: the same
posterior error produces an unboundedly larger velocity error, so up-weighting
it amplifies gradient noise rather than correcting a deficit, and it does so at
the expense of the low-$t$ region that actually determines the trajectory. The
problem is in the parameterisation, not in how the loss is spread over it ---
which is why no reweighting of this objective recovers the method, and why
changing what the network predicts (rather than how much each voxel counts) is
the direction the diagnosis points to.%
}

\newcommand{\REMEDYCONCLUSION}{%
Neither reweighting of the flow-matching objective recovers the method: the
decision-region reweighting is locally effective on the smallest class but costs
the primary metric, and the late-time reweighting degrades every region.%
}
